% 作者 shingja · github.com/shingja/shingja-math-modeling-workbench · 双许可见 LICENSE · 保留版权标识
% ============================================================
% 07 — 模型评价
% 说明: 客观评价模型优缺点, 提出改进方向。
%       分为"模型优点""模型缺点""模型改进与推广"三部分。
%       数值必须从 results/results.json 提取, 禁止手写。
% ============================================================

\section{模型评价与改进}

\subsection{模型优点}

本文建立的数学模型具有以下优点：

\begin{enumerate}
    \item \textbf{[优点1标题]:} [具体说明, 结合模型特点和数据结果]
    \item \textbf{[优点2标题]:} [具体说明, 结合模型特点和数据结果]
    \item \textbf{[优点3标题]:} [具体说明, 结合模型特点和数据结果]
    \item \textbf{[优点4标题]:} [具体说明, 结合模型特点和数据结果]
\end{enumerate}

\subsection{模型缺点}

本模型仍存在以下不足：

\begin{enumerate}
    \item \textbf{[缺点1标题]:} [具体说明, 客观分析]
    \item \textbf{[缺点2标题]:} [具体说明, 客观分析]
    \item \textbf{[缺点3标题]:} [具体说明, 客观分析]
\end{enumerate}

\subsection{模型改进与推广}

针对上述不足，未来可从以下方向进行改进：

\begin{enumerate}
    \item \textbf{算法改进:} [改进思路, 如采用更高效的算法/并行计算等]
    \item \textbf{模型扩展:} [扩展思路, 如考虑更多实际因素/放宽假设等]
    \item \textbf{应用推广:} [推广方向, 如将模型应用于其他领域/场景]
\end{enumerate}
